\documentclass[12pt]{article}
\usepackage[utf8]{inputenc}
\usepackage[T1]{fontenc}
\usepackage{hyperref}
\title{Y\'amana Language Resources for K-12 Curriculum Development}
\author{Compiled by Claude Sonnet 5.5 (claude-sonnet-5-5) for Andr\'as Kornai}
\date{}
\begin{document}
\maketitle

\section{Introduction}
Y\'amana (Yagan, Yahgan; ISO 639-3 yag) is an isolate of Tierra del Fuego in Chile and Argentina. Sources in the search results describe it as dormant after the death of Cristina Calder\'on on 16 February 2022, who was the last first-language speaker. Her granddaughter Cristina Z\'arraga and others are working on revitalization. No school curriculum was found.

\section{Primary Educational Resources}
No primers, textbooks or school programmes were found in the search results.

\section{Linguistic and Reference Materials}
\begin{itemize}
\item Wikipedia, Yahgan language. It mentions Thomas Bridges' dictionary of about 32,000 entries, published in 1933 and reprinted in 1965 as \emph{Yamana-English}. A fuller manuscript is at the British Library. Bridges also produced gospel translations and an 1894 grammar. \url{https://en.wikipedia.org/wiki/Yahgan_language}
\item Ben Molineaux (Edinburgh), work on Bridges' Yahgan materials: \url{https://benmolineaux.ppls.ed.ac.uk/Bridges-Yahgan-HiSoN.pdf}, \url{https://benmolineaux.ppls.ed.ac.uk/HiSoN23.pdf}, \url{https://benmolineaux.ppls.ed.ac.uk/ICHL27-solo.pdf}, \url{https://benmolineaux.ppls.ed.ac.uk/ICHL27-solo-slides.pdf}
\end{itemize}

\section{Religious and Community Resources}
Bridges' gospel translations are mentioned in the sources above, but the searches did not return a copy or a URL for them. Wikipedia on the Yaghan people: \url{https://en.wikipedia.com/wiki/Yaghans}

\section{Digital and Audio Resources}
Press features describe the language and its last speaker. They are not teaching resources. \url{https://nautil.us/excavating-a-language-at-the-end-of-the-world-747254}, \url{https://atlasoceanvoyages.com/blog/the-final-word}

\section{Curriculum Development Context}
Because the language is dormant, a curriculum would have to rebuild it from Bridges' materials, the recordings of Calder\'on and the work of the community. The sources found do not show what teaching materials exist. Contact with the Yagan community and with Molineaux's research would be the next step.
\end{document}